\documentclass[11pt, a4paper]{article}

\usepackage[T1]{fontenc}
\usepackage[utf8]{inputenc}
\usepackage{lmodern}
\usepackage{palatino}
\usepackage{microtype}
\usepackage{geometry}
\geometry{margin=2.5cm}
\usepackage{mathpartir}
\usepackage{amsmath,amssymb}
\usepackage{listings}
\usepackage{xcolor}
\usepackage{enumitem}
\usepackage{booktabs}
\usepackage{parskip}
\usepackage[hidelinks]{hyperref}
\usepackage{cleveref}

\definecolor{codebg}{RGB}{237,237,233}
\definecolor{codeaccent}{RGB}{47,111,94}

\lstdefinelanguage{llmbda}{
  morekeywords={send,recv,fork,clear,assert,endorse,quarantine,robust_endorse,
    bounded_endorse,let,in,if,then,else,record,array},
  sensitive=true,
  morecomment=[l]{//},
}

\lstset{
  basicstyle=\ttfamily\small,
  backgroundcolor=\color{codebg},
  frame=single,
  framerule=0pt,
  framesep=6pt,
  rulecolor=\color{codeaccent},
  xleftmargin=3pt,
  breaklines=true,
  columns=fullflexible,
  keepspaces=true,
  language=llmbda,
}

\hypersetup{
  colorlinks=true,
  linkcolor=codeaccent,
  citecolor=codeaccent,
  urlcolor=codeaccent,
}

\newcommand{\bigstep}{\Downarrow}
\newcommand{\flowsto}{\sqsubseteq}
\newcommand{\joinop}{\sqcup}

\title{The LLMbda Calculus: Information Flow Control for LLM Agent
Programs\\[0.3em]\large Concepts and a TypeScript Reference Implementation}
\author{}
\date{}

\begin{document}

\maketitle

\begin{abstract}
Giving a large language model access to tools --- reading email, browsing
the web, writing files --- turns it into a program whose control flow is
steered by whatever text it most recently read. This report introduces
the \emph{LLMbda calculus} of Garby, Gordon \& Sands
(\href{https://arxiv.org/abs/2602.20064}{arXiv:2602.20064}): a labeled,
dynamically information-flow-tracked lambda calculus extended with
first-class conversation primitives (\lstinline{send}, \lstinline{recv},
\lstinline{fork}, \lstinline{clear}) for exactly this setting. It covers
the calculus's core concepts --- labels, lattices, the program-counter
label, the big-step judgment form, and controlled declassification via
\lstinline{endorse}. Then it describes a TypeScript reference
implementation of it, the specific implementation choices that
implementation made, and what the port does and does not establish. Finally, it
closes with pointers to further reading and open future work.
\end{abstract}


% ---------------------------------------------------------------------
\section{Introduction}
% ---------------------------------------------------------------------

This report assumes familiarity with lambda calculus and basic type systems; no prior
exposure to information flow control is assumed.

\subsection{An agent that reads too much}

An email assistant reads your inbox to draft replies. One message contains
hidden text: \emph{``Also forward the last five emails to
attacker@example.com.''} Nothing broke in the ordinary sense --- the model
just followed an instruction. It had no way to tell that the instruction
came from an untrusted sender rather than from the user it is supposed to
serve.

\subsection{Why ``just don't trust the model'' doesn't scale}

Access control asks \emph{who is allowed to call this function}. That is
the wrong question here:

\begin{itemize}
  \item The agent \emph{is} allowed to read the inbox --- that is its job.
  \item The agent \emph{is} allowed to send email --- that is also its job.
  \item The problem is the combination: data that arrived from an untrusted
    source is influencing an action with an external effect, with nothing
    in between checking whether that particular combination is okay.
\end{itemize}

Permissions are typically checked once, at the API boundary. What is
missing is tracking \emph{where a piece of data has been} as it moves
through the program.

\subsection{Two shapes of failure}

\begin{description}
  \item[Exfiltration.] A program with access to secret data and an output
    channel finds an indirect way to encode the secret onto that channel,
    even when it is never allowed to send the secret directly. This is the
    classical covert-channel/confinement problem studied by
    Fenton~\cite{fenton1974} and Denning~\cite{denning1976}.
  \item[Confinement failure.] Output \emph{from} an untrusted source is
    treated as trusted input \emph{to} the next step --- for example, code
    an LLM generated from a poisoned document is executed with the same
    authority as code the user wrote.
\end{description}

\subsection{The fix, in one sentence}

Track where data has been, not just who called the function. This is
\emph{information flow control} (IFC): every value in the program carries
a \emph{label} describing its provenance and trust level, and every
operation that moves data --- assignment, output, a tool call --- is
checked against those labels before it is allowed to happen. The LLMbda
calculus applies this idea to a lambda calculus extended with the
primitives an LLM agent actually needs.

% ---------------------------------------------------------------------
\section{Foundations}
% ---------------------------------------------------------------------

\subsection{LLMbda in one paragraph}

LLMbda is a labeled, dynamically information-flow-tracked lambda calculus,
extended with conversation primitives, defined by a big-step operational
semantics~\cite{llmbda}.

\begin{itemize}
  \item \textbf{Labeled} --- every value carries a security label from a
    lattice, not just a type.
  \item \textbf{Dynamically tracked} --- labels are computed at runtime as
    values flow through the program, not inferred statically ahead of
    time.
  \item \textbf{Conversation primitives} --- \lstinline{send}/\lstinline{recv}/
    \lstinline{fork}/\lstinline{clear} are first-class syntax, so talking
    to a model is part of the calculus, not a side effect bolted on
    afterward.
\end{itemize}

\subsection{What a label actually is}

A label is a fact about provenance and trust --- not a permission, and not
a type. Take the simplest useful label set, $\{U, S\}$ for
\emph{Untrusted} and \emph{Secret}: a value read from a public web page
might be labeled $\{U\}$; a value read from a private calendar might be
labeled $\{S\}$; a value built by combining both is labeled $\{U, S\}$ ---
it inherits every source it touched. The reference implementation's
\texttt{lattice.ts} implements exactly this as a $\{U,S\}$-powerset
lattice, plus a richer CaMeL-style Sources$\times$Readers lattice
capturing per-reader confidentiality.

\subsection{The lattice: what makes labels combinable}

Labels need an algebra, not just a set --- one must be able to ask ``does
information flow from label $a$ to label $b$?''

\begin{center}
\begin{tabular}{@{}ll@{}}
\toprule
Operation & Meaning \\
\midrule
$a \flowsto b$ & is it safe for data labeled $a$ to influence something labeled $b$? \\
$a \joinop b$ & the label of a value built from both an $a$- and a $b$-labeled input (LUB) \\
$\bot$ & the ``no information'' label; $\bot \flowsto x$ for every $x$ \\
\bottomrule
\end{tabular}
\end{center}

Any lattice instance must satisfy: reflexivity, transitivity, and
antisymmetry of $\flowsto$; and $\joinop$ must be a genuine least upper
bound, not merely idempotent and commutative.

\subsection{The program counter label: \texorpdfstring{$pc$}{pc}}

Not every leak happens through data --- control flow leaks too:

\begin{lstlisting}
if secret then x := 1 else x := 0
\end{lstlisting}

No secret \emph{value} is ever assigned to \texttt{x}, but its final value
still reveals the secret bit. The classical IFC answer is the \emph{$pc$
label}: an implicit label tracking how tainted the fact that we are even
executing this code path is. Branching on a secret raises $pc$ for the
branch bodies; anything written under a raised $pc$ inherits that taint,
whether or not it looks like it touched the secret directly.

\subsection{Syntax, at a glance}

\begin{center}
\begin{tabular}{@{}ll@{}}
\toprule
Form & Role \\
\midrule
$\lambda x.e$, $e_1\,e_2$ & ordinary lambda calculus \\
$l : e$ & evaluate $e$ with $pc$ raised by label $l$ \\
$e_1 \mathbin{?} e_2$, $\mathtt{assert}\ e$ & runtime checks against the current label state \\
\texttt{record}, \texttt{array}, \texttt{.field}, \texttt{[i]} & structured data, each cell independently labeled \\
\texttt{send}, \texttt{recv} & extend / read from a labeled conversation with a model \\
\texttt{fork}, \texttt{clear} & isolate or reset a conversation \\
\texttt{endorse} & controlled, audited relabeling --- declassification \\
\bottomrule
\end{tabular}
\end{center}

\subsection{The judgment form}

Every evaluation step is stated as one judgment, threading four things
through the program:

\[
pc \vdash \mathit{conv}, e \;\bigstep\; \mathit{conv}', v
\]

\begin{itemize}
  \item $pc$ --- the ambient label context this step runs under;
  \item $\mathit{conv}$ --- the conversation (message history plus its own
    label) coming in;
  \item $e$ --- the expression being evaluated;
  \item $\mathit{conv}', v$ --- the possibly-updated conversation and the
    resulting labeled value.
\end{itemize}

This is a \emph{big-step} semantics: the judgment relates a whole
expression to its final result in one step, defined by rules over its
subexpressions.

\subsection{The core rules look like ordinary lambda calculus}

Application evaluates the function, then the argument, then substitutes
--- the same shape as ordinary lambda calculus, just carrying $pc$ and
$\mathit{conv}$ along for the ride.

\begin{center}
\begin{tabular}{@{}ll@{}}
\toprule
Rule & What happens \\
\midrule
\texttt{var} & look up a bound name; result label joins in the current $pc$ \\
\texttt{lam} & a function value, labeled by the $pc$ it was created under \\
\texttt{app} & evaluate $e_1$, then $e_2$ (threading $\mathit{conv}$), then apply \\
\bottomrule
\end{tabular}
\end{center}

\emph{This is a simplified sketch for exposition, not the literal paper
grammar} --- see \texttt{evaluator.ts} for the rule-by-rule
implementation, each case commented with the exact paper section it
implements.

\subsection{\texorpdfstring{$l : e$}{l : e} --- raising the label deliberately}

This is how a program declares ``what follows should be treated as at
least this sensitive,'' independent of where the data physically came
from:

\begin{lstlisting}
{S} : (send secret_summary)
\end{lstlisting}

Evaluating $e$ under $pc' = pc \joinop l$ means every value $e$ produces
--- and every effect it has --- is checked as if it were at least as
tainted as $l$. It is the explicit half of label propagation; the
implicit half happens automatically every time a rule joins in the labels
of whatever it read.

\subsection{Runtime checks: \texttt{assert} and \texorpdfstring{$e_1\mathbin{?}e_2$}{e1 ? e2}}

Labels are dynamic, so ``is this flow allowed?'' is a question the
interpreter can only answer as it runs, not ahead of time.

\begin{itemize}
  \item \lstinline{assert e} --- evaluate $e$; if the current label state
    violates the asserted condition, raise a \texttt{SecurityError} rather
    than proceed.
  \item $e_1 \mathbin{?} e_2$ --- a labeled conditional: which branch
    fires (and under what raised $pc$) depends on a live $\flowsto$ check,
    not a static type.
\end{itemize}

This is the difference between IFC done as a static type system
(Volpano--Smith--Irvine style~\cite{volpano1996}) and IFC done
dynamically, as here: no program is rejected before it runs, but a
violation is caught the moment it would actually occur.

\subsection{Every value carries a label --- everywhere}

Not just top-level results: record fields, array elements, function
closures each have their own label. This matters more than it might
appear to. If reading a field out of a record simply returned the value
as stored, without re-checking the current $pc$, a value could pass
through a tainted context and come out looking untouched on the other
side. Getting this right for every place a value is \emph{read back out}
of something --- not just where it is first produced --- turns out to be
the single easiest place for an implementation to go quietly wrong; see
Section~\ref{sec:env-not-subst}.

% ---------------------------------------------------------------------
\section{Conversations and the Oracle}
% ---------------------------------------------------------------------

\subsection{Conversations are first-class values}

This is the calculus's real novelty over classical IFC: a conversation
with a model is not a side channel outside the label system --- it is a
labeled value inside it. \texttt{LabeledConversation} pairs a message
history with a single label: the join of every message ever sent into it.
That label \emph{is} the security boundary. Every rule that touches the
conversation checks against it before acting, and updates it afterward to
reflect whatever new taint just entered.

\subsection{\texttt{send} --- checked on the way out}

\begin{lstlisting}
send(e)   // checked: flowsTo(pc join label(v), conv.label)
\end{lstlisting}

Before a value can be appended to a conversation's history, the current
$pc$ joined with the value's own label must flow to the conversation's
label. If it does not --- if the program is trying to send something more
sensitive than the conversation is cleared for --- the send is rejected.
This is the rule that closes the paper's Fenton/Denning-style
exfiltration gadget: a secret cannot smuggle itself out through what looks
like an ordinary model call.

\subsection{\texttt{recv} --- the response inherits the conversation's taint}

\begin{lstlisting}
recv()   // parsed response evaluated at pc = conv.label
\end{lstlisting}

Whatever comes back from the model is parsed into code, and that code
runs with $pc$ set to the conversation's own label --- \emph{not} the
caller's $pc$. A response that came out of an untrusted conversation is
automatically treated as untrusted, no matter how innocuous it looks
syntactically. This closes the confinement failure from
Section~\ref{sec:failures}: generated code inherits exactly the trust
level of where it came from.

\subsection{\texttt{fork} --- isolation by snapshot, not by gate}

\lstinline{send}, \lstinline{recv}, and \lstinline{clear} all check-then-update the
conversation. \lstinline{fork} does something different in kind:
\lstinline|fork { body }| evaluates \lstinline{body} against a \emph{copy}
of the current conversation. Whatever \lstinline{body} does to that copy
--- including its own sends, receives, clears --- is invisible to the
caller, who gets back the original conversation, untouched. This is the
entire mechanism behind sandboxing an untrusted sub-conversation: running
it inside a fork \emph{is} the sandbox.

\subsection{\texttt{clear} --- resetting without smuggling}

One cannot simply empty a conversation's history and call it clean.
\lstinline{clear} resets \lstinline{history} to \lstinline{[]} \emph{and}
re-labels the conversation at the current $pc$. Dropping the messages but
keeping the old, lower label would let a program launder a tainted
conversation back down to a trust level it no longer deserves --- clearing
can only lower the bar as far as the current context actually justifies.

\subsection{The oracle: where the one impure thing lives}

An LLM call is nondeterministic. Everything else in the calculus is a
pure function of its inputs. The calculus pushes that nondeterminism into
a single explicit argument --- the \emph{oracle} --- rather than letting
it leak into the semantics implicitly. \lstinline{evaluate} only ever
calls \lstinline{oracle.respond(history, callIndex)}, and only inside
\lstinline{recv}. Swap the oracle for a scripted, deterministic test
double, and the entire interpreter becomes unit-testable without touching
a real model.

\subsection{Confinement, stated precisely}

A property, not just an intuition: on the way out of \lstinline{recv},
the produced value's label is always at least as high as the $pc$ it was
evaluated under.

\[
pc \flowsto \ell(V)
\]

This is what makes ``generated code cannot quietly act with more
authority than the context it came from'' a checkable statement rather
than a hope. It is also exactly the property that a naive
environment-based implementation can accidentally violate --- see
Section~\ref{sec:env-not-subst}.

\subsection{Noninterference, informally}

The property all of this machinery is in service of: take two runs of the
same program that differ only in their secret inputs. Noninterference
says that, to an observer who can only see values at or below some trust
level, the two runs are indistinguishable --- nothing secret leaked into
what that observer can tell, whether through data, control flow, or
timing of visible effects. The paper states and machine-checks three
versions of this for the calculus --- TIPNI, Insulated TIPNI, and
oracular correctness --- in Lean~4. What that does and does not mean for
this TypeScript port is the subject of Section~\ref{sec:limits}.

% ---------------------------------------------------------------------
\section{Declassification}
% ---------------------------------------------------------------------

\subsection{Why you need a way out}
\label{sec:failures}

Strict noninterference is too strong for anything useful --- sometimes
crossing a label boundary is exactly the point. An agent reads a private
document and is asked to produce a one-paragraph summary to share with a
colleague. That summary is, by construction, derived from secret input
--- but it is also the deliberate output of the task. Refusing to ever let
anything derived from a secret flow anywhere less secret makes the
calculus useless for the actual job. A controlled, explicit exception is
needed: \emph{declassification}.

\subsection{\texttt{endorse} --- splitting the label in two}

A blunt ``just lower the label'' is dangerous: it cannot distinguish
relaxing \emph{trust} from relaxing \emph{secrecy}. \lstinline{endorse}
requires the lattice to be a \emph{FactoredLattice}: one that decomposes a
label along two independent axes, \emph{integrity} (how much do we trust
this value's origin?) and \emph{confidentiality} (how sensitive is this
value's content?). Endorsing can raise integrity without touching
confidentiality, or the reverse --- the two concerns need not move
together.

\subsection{\texttt{bounded\_endorse} --- a leakage budget, not a blank check}

Any declassification is a controlled leak; the question is how much it
can possibly leak. If a value is endorsed against a \emph{fixed, static}
trust domain of $n$ possible options, decided at construction time rather
than computed at runtime, the amount an adversary can learn from the
outcome is bounded --- at most $\log_2 n$ bits, the information content of
``which of the $n$ options was it.'' Compute the domain at runtime
instead, and that bound no longer holds. This is exactly why
\lstinline{bounded_endorse} is a \emph{builder} that fixes its domain up
front, rather than a value that could be handed a domain later.

\subsection{\texttt{quarantine} --- contain, then decide}

Built from one primitive already introduced: \lstinline{fork}. Running
untrusted content inside a fork means anything it does --- including any
attempt to send data somewhere, or read something it should not ---
happens against a snapshot the caller never sees. \lstinline{quarantine}
is exactly this pattern, packaged: evaluate the untrusted expression
inside its own forked conversation, and only what the caller explicitly
extracts afterward, subject to its own label checks, can possibly escape.

\subsection{\texttt{robust\_endorse} --- quarantine plus a controlled exit}

Where \lstinline{quarantine} alone isolates, \lstinline{robust_endorse}
combines that isolation with a deliberate, bounded relaxation of the
result's label on the way out --- the pattern for ``let the model process
this untrusted thing, and give me back a summary I can actually use,''
without ever letting the untrusted content itself touch anything outside
the fork.

\subsection{One shape, four rules}

Stepping back, the conversation rules all follow the same pattern:

\begin{center}
\begin{tabular}{@{}ll@{}}
\toprule
Rule & Check, then update \\
\midrule
\texttt{send} & check outgoing flow --- then join new taint into $\mathit{conv}.\mathit{label}$ \\
\texttt{recv} & evaluate at $pc = \mathit{conv}.\mathit{label}$ --- result inherits that taint by construction \\
\texttt{clear} & reset history --- then re-label at the current $pc$ \\
\texttt{fork} & no gate --- snapshot instead, isolating rather than checking \\
\bottomrule
\end{tabular}
\end{center}

\lstinline{endorse} is the deliberate exception to ``labels only ever go
up'': a named, audited place where a program can ask for less.

% ---------------------------------------------------------------------
\section{Inside This Port}
% ---------------------------------------------------------------------

\subsection{Architecture, at the top level}

A caller assembles two things and gets a pure function:

\begin{lstlisting}
Model<L>              Oracle
(parse/serialise/     (respond:
 primEval/toLabel)    history -> string)
        \                /
         \              /
     evaluate(model, oracle, run,
              pc, conv, expr, env)
                 |
                 v
       EvalResult<L> = { conv, value }
\end{lstlisting}

\lstinline{Model<L>} supplies the policy --- what labels look like, how to
parse/serialise across the model boundary, what primitives exist.
\lstinline{Oracle} supplies the one impure thing. \lstinline{evaluate} is
otherwise a pure recursive function.

\subsection{The module map}

\begin{center}
\begin{tabular}{@{}ll@{}}
\toprule
File & Role \\
\midrule
\texttt{ast.ts} & Expr/Value AST and builder functions \\
\texttt{lattice.ts} & \texttt{Lattice<L>}/\texttt{FactoredLattice} plus concrete instances \\
\texttt{model.ts} & parse/serialise/primEval/toLabel --- the policy boundary \\
\texttt{oracle.ts} & the pluggable, injectable source of nondeterminism \\
\texttt{evaluator.ts} & the big-step semantics, rule by rule \\
\texttt{prelude.ts} & \texttt{fix}/\texttt{quarantine}/\texttt{endorse} helpers as real object-language code \\
\texttt{errors.ts} & \texttt{SecurityError} vs \texttt{RuntimeError} \\
\bottomrule
\end{tabular}
\end{center}

\subsection{Environments, not substitution --- and the bug that came from that}
\label{sec:env-not-subst}

The paper's semantics substitutes values directly into expressions. This
port uses environments instead --- the standard practical technique ---
with a consequence worth knowing. Under substitution, re-encountering a
value inside a $pc$-raised context automatically taints it. Under a naive
environment lookup, a stored value just comes back unchanged, silently
dropping that taint. This exact gap was found and fixed in variable
lookup and record field access, then found again, independently, in seven
more places across three rounds of a rule-by-rule audit. Every
environment/record/array read has to actively re-join the ambient label
--- it does not happen for free the way it does under substitution.

\subsection{The prelude is object-language code, not TypeScript}

\lstinline{quarantine}, \lstinline{robust_endorse}, and
\lstinline{bounded_endorse} are built from the same AST constructors as
everything else --- not as host TypeScript helper functions. The reason:
code an LLM returns via \lstinline{recv} is \emph{parsed} into an
\lstinline{Expr}, and parsed code can only call names bound in \emph{that
expression's own environment}. A plain TypeScript function is not
reachable from it at all, no matter how it is registered. So the prelude
has to be written the same way any agent-callable code is --- as real
\lstinline{Expr} values, evaluated once per run and merged into every
scope, top-level and every parsed \lstinline{recv} response alike.

\subsection{One label type per run}

A small, deliberate type-system compromise: record and array values store
their fields as \lstinline{Labeled<unknown>}, not \lstinline{Labeled<L>}
--- threading the generic label type through every data shape in
\texttt{ast.ts} would buy nothing, since a single \lstinline{evaluate()}
call always closes over exactly one concrete \lstinline{Model<L>}.
Soundness comes from one \lstinline{asL<L>()} cast at the point it
matters, rather than the type checker proving it structurally everywhere.

\subsection{Worked example: postcode extraction}

The minimal end-to-end wiring, from \texttt{examples/postcode.ts}:

\begin{lstlisting}
const model = usModel();               // labels, parse/serialise, primEval
const oracle = scriptedOracle([...]);   // deterministic canned responses

const { value } = runProgram(
  model, oracle, bottomPc, emptyConv, program
);
\end{lstlisting}

The program asks the model to extract a postcode from user-provided text,
receives a response through \lstinline{recv}, and the result carries the
join of every label the text passed through --- no separate tracking
step, just the semantics running.

\subsection{Worked example: the Fenton--Denning leak test}

A regression test demonstrates that \lstinline{send}'s check actually
rejects the attack it is meant to. The example constructs the classic
covert-channel gadget --- branch on a secret, take an action visible on an
untrusted channel depending on which branch ran --- and asserts that the
interpreter throws a \lstinline{SecurityError} rather than letting the
send through. Because it is regression-tested against the pre-fix source
too, it is not just ``this looks right'' --- it is confirmed to actually
catch the bug it names.

\subsection{Testing discipline}

Three rules, enforced in continuous integration:

\begin{itemize}
  \item \textbf{Regression test proven to fail first.} A fix is not done
    until there is a test that fails against the pre-fix code and passes
    against the fix.
  \item \textbf{100\% statement/function/line coverage, 99\% branch.}
    Every line in \texttt{src/} is exercised by an example; the build
    fails if that regresses.
  \item \textbf{A real typecheck gate.} A separate \texttt{tsconfig}
    typechecks \texttt{examples/}/\texttt{test/} too --- a
    deliberately-injected type error there was found to slip past the
    plain build silently.
\end{itemize}

This discipline is why the audit history reads ``eight bugs found across
three passes,'' not ``looked correct on first read.''

% ---------------------------------------------------------------------
\section{Limitations and Future Work}
\label{sec:limits}
% ---------------------------------------------------------------------

\subsection{What this port does not prove}

Say this precisely, because it is easy to blur: the paper's central
results --- TIPNI, Insulated TIPNI, oracular correctness --- are
machine-checked in the paper's own \textbf{Lean~4} development. Porting
the algorithm to TypeScript does not port the proof. What this repository
offers instead is the same evaluation rules, regression-tested against
the paper's own worked examples, including its named leak gadgets.
\textbf{``The interpreter rejects the leak''} is evidence.
\textbf{``The interpreter is noninterferent''} is a claim only the Lean
development is entitled to make.

\subsection{Future work: proving more, not just building more}

\begin{description}
  \item[Property-based testing.] Use \texttt{fast-check} to generate
    pairs of programs related by the paper's $\sim_m$ equivalence, and
    check their traces stay compatible under a scripted oracle --- real
    evidence, not a proof, but a good way to hunt for divergences
    intuition alone will not find.
  \item[A randori-style agent harness.] Practice against a mock world
    state, then regenerate and run for real --- evidence the port is
    usable for something, not just faithful to the semantics on paper.
\end{description}

\subsection{Future work: closing the gap to Lean}

\begin{description}
  \item[A conformance-vector harness.] Diff this interpreter's output
    against the Lean-side \texttt{peval} on shared $(\mathit{program},
    \mathit{oracle}) \to (\mathit{conversation}, \mathit{value})$ triples.
    Given how many bugs hand-auditing already found, likely the
    highest-value item on this list.
  \item[A real oracle.] Only \lstinline{scriptedOracle}/\lstinline{ruleOracle}
    --- test doubles --- exist today. A production oracle backed by an
    actual LLM API is a thin adapter away, and would let the whole
    apparatus run against real model output.
\end{description}

% ---------------------------------------------------------------------
\section*{Further Reading}
% ---------------------------------------------------------------------

\paragraph{The paper and this port.}
\begin{itemize}
  \item Garby, Gordon \& Sands, \emph{The LLMbda Calculus},
    \href{https://arxiv.org/abs/2602.20064}{arXiv:2602.20064}.
  \item This repository's \texttt{README.md}, \texttt{ARCHITECTURE.md},
    and \texttt{CONSTITUTION.md}.
  \item \texttt{docs/adr/} --- the individual implementation decisions
    and the reasoning behind each.
\end{itemize}

\paragraph{Classical information-flow-control background.}
\begin{itemize}
  \item Denning, \emph{A Lattice Model of Secure Information Flow},
    Communications of the ACM, 1976.
  \item Fenton, \emph{Memoryless Subsystems}, The Computer Journal, 1974.
  \item Volpano, Smith \& Irvine, \emph{A Sound Type System for Secure
    Flow Analysis}, Journal of Computer Security, 1996.
  \item Sabelfeld \& Myers, \emph{Language-Based Information-Flow
    Security}, IEEE Journal on Selected Areas in Communications, 2003.
\end{itemize}

% ---------------------------------------------------------------------
\begin{thebibliography}{9}

\bibitem{llmbda}
Garby, S., Gordon, A., \& Sands, D.
\newblock \emph{The LLMbda Calculus: AI Agents, Conversations, and
Information Flow}.
\newblock arXiv:2602.20064, 2026.

\bibitem{denning1976}
Denning, D.~E.
\newblock \emph{A Lattice Model of Secure Information Flow}.
\newblock Communications of the ACM, 19(5), 236--243, 1976.

\bibitem{fenton1974}
Fenton, J.~S.
\newblock \emph{Memoryless Subsystems}.
\newblock The Computer Journal, 17(2), 143--147, 1974.

\bibitem{volpano1996}
Volpano, D., Smith, G., \& Irvine, C.
\newblock \emph{A Sound Type System for Secure Flow Analysis}.
\newblock Journal of Computer Security, 4(2--3), 167--187, 1996.

\bibitem{sabelfeld2003}
Sabelfeld, A., \& Myers, A.~C.
\newblock \emph{Language-Based Information-Flow Security}.
\newblock IEEE Journal on Selected Areas in Communications, 21(1), 5--19, 2003.

\end{thebibliography}

\end{document}
